\chapter*{Zusammenfassung}

Diese Arbeit untersucht MULTIFIT, ein Approximationsverfahren von Coffman, Garey und Johnson
(1978) für die Makespan-Minimierung auf identischen Maschinen.
Das Verfahren rät eine Schranke an den Makespan und überlässt dem Bin-Packing-Verfahren First
Fit Decreasing (FFD) als Subroutine die Frage, ob die Aufgaben in ebenso viele Behälter dieser
Größe passen, wie Maschinen vorhanden sind; eine Binärsuche verengt die Schranke anschließend.
Im ersten Teil der Arbeit wird die Performanzanalyse von MULTIFIT geschlossen aufbereitet: Die
auf Coffman, Garey und Johnson (1978) zurückgehende Schranke $5/4$ für das Grenzverhältnis, an
welchem das im schlechtesten Fall erreichte Verhältnis von erzeugtem zu optimalem Makespan
hängt, wird vollständig bewiesen und mit dem Makespan nach $k$ Suchschritten verbunden; die
schärfere Schranke $13/11$ von Cao (1995) wird eingeordnet, und an einer Instanz von Friesen
(1984) wird gezeigt, dass sie sich für mindestens $13$ Maschinen nicht verbessern lässt.
Im zweiten Teil der Arbeit wird die Subroutine von MULTIFIT ausgetauscht.
Für Modified First Fit Decreasing (MFFD) von Johnson und Garey (1985) als Subroutine wird
gezeigt, dass es innerhalb von MULTIFIT dieselbe Zuordnung wie FFD liefert, sobald die mittlere
Maschinenlast mindestens doppelt so groß ist wie die größte Aufgabe.
Der $\mathrm{OPT}+1$-Algorithmus von Jansen und Solis-Oba (2010) für Instanzen mit $d$
verschiedenen Aufgabengrößen wird zu einem $\varepsilon$-dualen Verfahren im Sinne von Hochbaum
und Shmoys (1987) umgebaut, welches höchstens die optimale Binzahl verwendet und dafür die
Kapazität um den Faktor $1+\varepsilon$ überschreiten darf; MULTIFIT erreicht damit einen
Makespan von höchstens $(1+\varepsilon)(1+2^{-k})$ mal dem Optimum, wobei
$\varepsilon = 1/(2d-1)$ ist.
Für das Konfigurations-LP von Gilmore und Gomory (1961) trägt die beim Runden entstehende
Schranke dieselbe Schlussweise dagegen nicht.
Sämtliche untersuchten Verfahren sind implementiert und praktisch getestet worden.
Auf \EvalInstanzen{} Benchmark-Instanzen überschreitet kein gemessenes Verhältnis zum optimalen
Makespan den Wert $13/11$, und MFFD liefert keinen systematisch kleineren Makespan als FFD.
Auf \HmInstanzen{} Instanzen mit wenigen verschiedenen Aufgabengrößen liefert MULTIFIT mit dem
$\mathrm{OPT}+1$-Algorithmus auf \HmJansenBesserFFD{} Instanzen einen kleineren, auf
\HmJansenSchlechterFFD{} jedoch einen größeren Makespan als mit FFD.
Eine bessere Bin-Packing-Garantie sichert somit keinen kleineren Makespan; dafür kommt es darauf
an, bei welchen der getesteten Schranken die Subroutine die Maschinenzahl unterschreitet und, wo
sie die Schranke überschreiten darf, wie weit sie das tut.
